\documentclass[10pt]{article}
\usepackage[margin=1.4cm]{geometry}
\usepackage[T1]{fontenc}
\usepackage{booktabs,longtable,array,tabularx,multirow}
\usepackage{pdflscape}
\usepackage{enumitem}
\usepackage{amsmath,amssymb}
\usepackage[expansion=false]{microtype}
\usepackage[hidelinks]{hyperref}
\usepackage{xurl}
\newcolumntype{P}[1]{>{\raggedright\arraybackslash}p{#1}}
\setlength{\LTpre}{0.4em}
\setlength{\LTpost}{0.4em}
\setlist[itemize]{leftmargin=1.35em}
\setlist[enumerate]{leftmargin=1.45em}
\title{Supplementary Data S2: Screening Exclusion Log and Decision Codebook}
\author{Supplement to: \textit{Sixty Years Retrospective on Metaheuristic Algorithms (1965--2025): A Systematic Review}}


\begin{document}
\maketitle
\section*{Purpose}
The manuscript reports 4,696 identified records, 3,810 duplicate or duplicate-version instances removed, 886 deduplicated reports screened, 62 studies retained in the contemporary critical corpus, and 824 reports not retained for that deep-analysis corpus. The source package supplied with the manuscript does \emph{not} contain the bibliographic identities and record-level exclusion decisions for those 824 non-retained reports. This supplement therefore records the exact available aggregate evidence and the complete decision codebook without fabricating an item-level ledger.

\begin{table}[h]
\centering
\caption{Auditable study-selection arithmetic.}
\begin{tabular}{P{8.0cm}r}
\toprule
Stage & Count\\
\midrule
Records identified & 4,696\\
Duplicate/duplicate-version instances removed & 3,810\\
Deduplicated reports screened & 886\\
Reports not retained in the deep critical corpus & 824\\
Studies retained in the contemporary critical corpus & 62\\
\bottomrule
\end{tabular}
\end{table}

\section*{Primary exclusion codes}
\begin{description}
\item[E1] No first-level algorithm or substantive methodological contribution.
\item[E2] Application-only use of an established algorithm.
\item[E3] Duplicate/companion report or nominal variant without an independent contribution.
\item[E4] Non-peer-reviewed or preprint-only at the eligibility cutoff.
\item[E5] Insufficient mathematical, algorithmic or implementation detail.
\item[E6] Insufficient benchmark or real-world validation for the relevant research question.
\item[E7] Full text unavailable.
\item[E8] Outside the temporal or language scope.
\item[E9] Insufficient methodological quality for entry into the deep critical corpus under the stated 16/20 threshold.
\end{description}

\section*{Decision hierarchy}
When several reasons apply, the earliest decisive reason is the primary code. Any additional applicable reason is a secondary code. The exclusion decision must be linked to a \texttt{REC} or \texttt{REP} identifier, screening stage, reviewer decision and source record. No distribution of E1--E9 is reported here because the record-level decisions are absent from the supplied archive.

\begin{longtable}{P{4.0cm}P{10.2cm}}
\caption{Completeness status of the exclusion evidence.}\\
\toprule
Evidence item & Status\\
\midrule
\endfirsthead
\toprule Evidence item & Status\\\midrule\endhead
Aggregate excluded count & Verified from manuscript: 824 reports.\\
Exclusion-code definitions & Fully specified as E1--E9 above.\\
Item-level report identities & Not recoverable from the supplied manuscript/source archive.\\
Item-level primary reasons & Not recoverable from the supplied manuscript/source archive.\\
Item-level secondary reasons & Not recoverable from the supplied manuscript/source archive.\\
Quality score for excluded reports & Not recoverable from the supplied manuscript/source archive.\\
Retrospective reconstruction & Prohibited in this supplement; no pseudo-records or inferred reasons are created.\\
\bottomrule
\end{longtable}

\paragraph{Submission interpretation.}
This file is a complete disclosure of the exclusion evidence that is supportable from the manuscript. A true 824-row exclusion ledger can only be produced from the authors' original screening/export records; inventing such rows would contradict the manuscript's stated non-fabrication rule.
\end{document}
